\section{Equivalence with Monoidal Frontier BSP}

Turing completeness is not the primary target.  It can often be obtained by
encoding a tape in a path graph, but says little about bulk parallel structure,
device residency, or asymptotic work.  We instead seek a characterization that
matches the intended execution model.

\begin{definition}[Monoidal frontier BSP source machine]
A monoidal frontier bulk-synchronous program operates on a finite ordered
multigraph.  Let $S$, $P$, and $A$ be the vertex-state, edge-payload, and
message types.  Its semantic operations are an edge message function
\[
 \mu:V\times V\times S\times S\times P\to A,
\]
a lawful commutative monoid $(A,\oplus,0)$, a vertex update
$u:V\times S\times A\to S$, and a next-frontier selector
$q:V^*\times(V\to S)\times(V\to S)\to V^*$.  Its first argument is the
ordered stream of destination candidates encountered by scanning the active
adjacency rows.  A barrier configuration is a pair $(\sigma,F)$ of a dense
store and an ordered frontier.
\end{definition}

For destination $v$, the source machine constructs $I_{\sigma,F}(v)$ by a
nested left fold: first over sources in $F$, then over each source adjacency
row, combining $\mu(s,d,\sigma(s),\sigma(d),p)$ exactly when $d=v$.  It then
sets
\[
 \sigma'(v)=u(v,\sigma(v),I_{\sigma,F}(v)),\qquad
 F'=q(D_G(F),\sigma,\sigma'),
\]
where $D_G(F)$ is built by the same nested frontier/adjacency scan and retains
every destination occurrence in exact scan order.
This definition does not invoke the flattened traversal or any target
operator.

\begin{definition}[Compiled pull--map--push step]
At state $\sigma$, compile the edge map
\[
 \widehat\mu_\sigma(c)=
 \mu(\mathsf{src}(c),\mathsf{dst}(c),
     \sigma(\mathsf{src}(c)),\sigma(\mathsf{dst}(c)),
     \mathsf{payload}(c)).
\]
The target flattens the frontier and adjacency rows into $\trav_G(F)$, maps
$\widehat\mu_\sigma$, pushes messages by destination using $\oplus$, and uses
the same $u$ and $q$ at the barrier.  Its candidate stream is the destination
projection of the flattened traversal.  Pulls are re-instantiated against the new
store after every superstep at the semantic level; this does not require
runtime kernel compilation.
\end{definition}

\begin{lemma}[Nested-fold flattening]\label{lem:nested-fold-flattening}
For every graph, state, frontier, destination, and initial accumulator,
folding the flattened traversal gives the same accumulator as the source
machine's nested frontier/adjacency fold.
\end{lemma}

\begin{proof}
Induct on each adjacency row to show that mapping it to edge contexts preserves
the destination-filtered fold.  Then induct on the frontier and use left-fold
composition over sequence concatenation.  The Lean declarations are
\texttt{fold\_expand\_compile} and \texttt{fold\_traverse\_compile}.
\end{proof}

\begin{theorem}[One-step forward simulation]\label{thm:step-simulation}
For every finite ordered graph, monoidal frontier BSP program, and barrier
configuration, the source superstep and its compiled pull--map--push step
produce equal stores and equal ordered frontiers.
\end{theorem}

\begin{proof}
Lemma~\ref{lem:nested-fold-flattening}, instantiated with the monoid identity,
gives pointwise equality of every destination inbox.  The source's nested
candidate scan and the target's destination projection are also equal as exact
lists by \path{traversalDestinations_eq_nestedDestinations}.  The update and
selection functions therefore receive extensionally equal arguments.  This is
checked by \texttt{compile\_step\_correct} in Lean.
\end{proof}

\begin{corollary}[Finite-execution forward simulation]
\label{cor:run-simulation}
For every $k\in\mathbb{N}$ and initial configuration, $k$ source supersteps
equal $k$ compiled target supersteps.
\end{corollary}

\begin{proof}
Induction on $k$ using Theorem~\ref{thm:step-simulation}; checked by
\texttt{compile\_run\_correct}.
\end{proof}

\begin{theorem}[Linear-schedule independence]
\label{thm:schedule-independence}
Let $L$ be any permutation of $\trav_G(F)$.  For every destination $v$, the
destination-filtered fold over $L$ equals the ordered target inbox and the
source-machine inbox.
\end{theorem}

\begin{proof}
Two destination-filtered fold actions commute.  If neither or only one context
targets $v$, this is immediate.  If both do, associativity moves the common
accumulator outside and commutativity exchanges their messages.  Induction on
the list-permutation derivation then preserves the complete fold.  Lean checks
the generic result in \path{foldContextAt_perm}, the target result in
\path{inboxAt_schedule_independent}, and its connection to the source in
\path{inbox_schedule_compile_correct}.
\end{proof}

These results quantify over arbitrary Lean types and total semantic functions;
they are not graph enumeration or bounded testing.  Exact ordered folds make
the present simulation stronger than equality up to a bag, and do not require
commutativity inside the forward-simulation proof.  Theorem~\ref{thm:schedule-independence}
uses the monoid laws to permit every linear permutation.  A refinement theorem
for arbitrary abstract reduction-tree shapes remains separate mathematical
work; no concrete GPU execution model is required here.

The semantic simulation deliberately quantifies over arbitrary total Lean
functions.  This makes it a broad specification envelope, but that envelope is
too broad to be a GPU language: an arbitrary selector $q$ may perform any
global computation.  We therefore refine it with syntax rather than claiming
that all functions are compilable.

\subsection{Signature-relative typed refinement}

\begin{definition}[Typed scalar IR]
For a fixed signature $\Sigma$, object-language types and terms are
\[
 \tau ::= \mathsf{Bool}\mid\mathsf{Index}_n\mid b
          \mid \tau_1\times\tau_2,
 \qquad
 t ::= x\mid\mathsf{literal}(\ell)\mid(t_1,t_2)
          \mid\mathsf{primitive}(o,t)
          \mid\mathsf{let}\ x=t_1\ \mathsf{in}\ t_2.
\]
Here $b$, $\ell$, and $o$ are symbols declared by $\Sigma$.  The signature
assigns denotations to base types, literal symbols, unary primitive symbols,
and reduction-monoid symbols; every declared monoid carries its associative,
identity, and commutativity proofs.  Multi-argument primitives consume product
types, matching the algebra's zip-then-map shape.  A let binding evaluates its
input once and exposes it to a singleton typed environment; it is the explicit
sharing node used by normalization.  Terms contain no denotation functions
directly: they refer to signature symbols instead.
\end{definition}

Variables are intrinsically typed de Bruijn references into heterogeneous
environments.  The message, update, and selection contexts are respectively
\[
 \Gamma_m=[\mathsf{Index},\mathsf{Index},S,S,P],\quad
 \Gamma_u=[\mathsf{Index},S,A],\quad
 \Gamma_q=[\mathsf{Index},S,S].
\]
Consequently, an ill-typed pull or primitive application cannot be
constructed.

\begin{definition}[Closed typed BSP fragment]
A typed source program consists of a message term
$m:\Gamma_m\vdash A$, a declared commutative monoid on $A$, an update term
$u:\Gamma_u\vdash S$, and a closed frontier policy.  The policy
$\mathsf{dense}(p)$, for $p:\Gamma_q\vdash\mathsf{Bool}$, enumerates every
valid vertex in ascending order and compacts those satisfying $p$;
$\mathsf{sparsePreserve}(p)$ instead stably filters the ordered destination
candidate stream, preserving the multiplicity of every accepted occurrence.
Its target program contains the same typed terms under explicit traverse,
edge-map,
destination-push, vertex-update, and next-frontier fields.
\end{definition}

Dense selection remains the canonical $|V|$ baseline.  Sparse selection makes
candidate generation explicit instead of hiding a global function in the
model; its exact order, multiplicity, and cost properties are stated below.

\begin{theorem}[Compilation commutes with denotation]
\label{thm:typed-denotation}
For every signature $\Sigma$, well-typed source program $P$, and barrier state
$\sigma$,
\[
 \mathsf{denoteAt}(\mathsf{compile}_{\mathrm{syn}}(P),\sigma)
 =\mathsf{compile}_{\mathrm{sem}}(\mathsf{denote}(P),\sigma).
\]
\end{theorem}

\begin{proof}
Compilation preserves each typed scalar term and changes only graph-control
structure.  Evaluating the message term in its five-field heterogeneous
environment is definitionally the semantic edge map; reduction, update, and
frontier fields have the same denotations.  Lean checks the equality as
\path{Typed.compile_denoteAt}.
\end{proof}

\begin{corollary}[Typed finite-run compiler correctness]
\label{cor:typed-run}
For every finite ordered graph, typed source program, initial configuration,
and $k\in\mathbb{N}$, interpreting $k$ source supersteps equals interpreting
$k$ supersteps of its compiled typed pull--map--push target.
\end{corollary}

\begin{proof}
Combine Theorem~\ref{thm:typed-denotation} with the semantic one-step
simulation, then induct on $k$.  The declarations are
\path{Typed.compile_step_correct} and \path{Typed.compile_run_correct}.
The typed schedule result \path{Typed.inbox_schedule_correct} additionally
inherits Theorem~\ref{thm:schedule-independence}.  A concrete natural-number
path-mass program and evaluator equations are checked in
\path{TraversalAlgebra.TypedExample}, so the fragment is inhabited.
\end{proof}

\subsection{Closed Traversal Algebra fragment}

\begin{definition}[Compositional traversal expressions]
Fix one signature $\Sigma$ and types $S,P,A$.  The edge-stream expression
grammar is
\[
 e ::= \mathsf{read}(x)\mid\mathsf{literal}(\ell)
       \mid\mathsf{map}(e,t)\mid\mathsf{zip}(e_1,e_2).
\]
The intrinsically typed variable $x$ selects exactly one of source identifier,
destination identifier, source state, destination state, or edge payload.
In $\mathsf{map}(e,t)$, if $e$ has element type $X$, then
$t:[X]\vdash Y$ is a unary typed scalar term.  Zip constructs a product type.
There is no arbitrary edge callback and no pre-normalized message term in this
grammar.

A closed frontier-to-frontier expression has the nested operator form
\[
 Q=\mathsf{compact}_p
      (\mathsf{updateDst}_u(\reduceDst_M(e))).
\]
Here $e$ has element type $A$, $M$ is a declared commutative monoid on $A$,
$u:\Gamma_u\vdash S$, and the compact policy is either
$\mathsf{dense}(p)$ or $\mathsf{sparsePreserve}(p)$ for
$p:\Gamma_q\vdash\mathsf{Bool}$.  Its denotation is defined directly:
recursively evaluate $e$ at every context in $\trav_G(F)$, call the type-safe
\TA{} destination terminal, apply $u$ over the vertex store, and execute the
selected compact policy.  This operational definition invokes neither BSP nor
the normalizer.  Lean defines the grammar and semantics as
\path{Typed.TraversalAlgebra.Expression.Traversal},
\path{Destination}, and \path{Program.step}.
\end{definition}

This fragment is deliberately closed under the observation of a BSP barrier.
The \emph{emit} terminal returns an edge stream, and source reduction returns
one item per frontier occurrence; neither is a dense-store transition without
an additional output observation.  They are therefore not silently conflated
with the barrier theorem.  Section~\ref{sec:terminal-observations} treats them
with explicit result constructors and independent BSP observers.

\begin{definition}[Single-push normal form]
A normal-form \TA{} superstep is
\[
 \trav_G(F)
 \mathbin{\longrightarrow}^{\textnormal{one typed edge term }m} A^*
 \mathbin{\longrightarrow}^{\reduceDst_M}(V\to A)
 \mathbin{\longrightarrow}^{u}(V\to S)
 \mathbin{\longrightarrow}^{\mathsf{compact}_p}V^*.
\]
It is represented by \path{Typed.TraversalAlgebra.Program}.  This is the
four-component normal form used by the earlier BSP correspondence, rather
than the definition of arbitrary \TA{} syntax.
\end{definition}

Define normalization $N$ recursively on traversal expressions:
\[
\begin{aligned}
 N(\mathsf{read}(x)) &= x,\\
 N(\mathsf{literal}(\ell)) &= \mathsf{literal}(\ell),\\
 N(\mathsf{map}(e,t)) &= \mathsf{let}\ x=N(e)\ \mathsf{in}\ t,\\
 N(\mathsf{zip}(e_1,e_2)) &= (N(e_1),N(e_2)).
\end{aligned}
\]
The binding is \path{Typed.Term.share}, represented by the \path{Term.let1}
constructor.  Unlike tree substitution, it never copies $N(e)$ when $t$ reads
$x$ repeatedly: the input is evaluated once, and every use refers to the
bound value.  Normalization preserves the destination monoid, vertex update,
and frontier policy.

\begin{lemma}[Pointwise fusion]
\label{lem:ta-pointwise-normalization}
For every compositional traversal expression $e$, store $\sigma$, and edge
context $c$,
\[
 \denote{N(e)}(\rho_{\sigma,c})=\denote{e}_{\sigma}(c).
\]
Consequently the complete ordered edge streams before and after normalization
are equal.
\end{lemma}

\begin{proof}
Structural induction on $e$.  Read and literal cases are definitional.  Zip
uses both induction hypotheses.  The map case uses the shared-binding
equation \path{Term.evaluate_share}.  Lean checks the pointwise and
whole-stream results as \path{Traversal.evaluateAt_normalize} and
\path{Traversal.evaluate_normalize}.
\end{proof}

\begin{theorem}[Closed-expression normalization]
\label{thm:ta-normalization}
For every closed compositional \TA{} expression $Q$, graph $G$, barrier
configuration $C$, and $k\in\mathbb{N}$,
\[
 \mathsf{run}_{\mathrm{TAExpr}}(Q,G,k,C)
 =\mathsf{run}_{\mathrm{TANF}}(N(Q),G,k,C).
\]
The equality includes the dense store and ordered frontier.
\end{theorem}

\begin{proof}
Lemma~\ref{lem:ta-pointwise-normalization} makes the mapper functions
extensionally equal.  Applying the same destination fold gives equality of
every inbox; Lean checks this as \path{Destination.inboxAt_normalize}.  The
update and compact operations are unchanged, giving
\path{Program.normalize_step_correct}.  Induction on $k$ gives
\path{Program.normalize_run_correct}.

Conversely, reification $R$ recursively turns normal-form variables into
structural reads, pairs into zip nodes, and primitive applications and shared
bindings into map nodes.  Introducing an explicit binding means that even
$N(R(P))=P$ need not hold syntactically.  Lean instead proves pointwise
reification as \path{Traversal.evaluateAt_ofTerm} and
\path{Traversal.evaluate_normalize_ofTerm}, then lifts it to steps and finite
runs as \path{Program.ofNormalForm_run_correct} and
\path{reify_normalize_run_correct}.  Both reification round trips are thus
semantic, which is the relevant equality for language equivalence.
\end{proof}

\subsection{Sharing-preserving quantitative normalization}

The semantic theorem alone would permit an exponentially larger target term.
To exclude that failure mode, the cost model counts every expression and term
constructor, charges unit work to reads, literals, pair construction, and
primitive invocation, permits the two branches of a pair to run in parallel,
and counts a conservative peak of live per-edge scalar temporaries.  A
\path{let1} node counts its input and body once, evaluates them sequentially,
and never copies the bound syntax.  At program level, if $m$ is the number of
active edge occurrences, work additionally charges one monoid action per
edge, one update per vertex, and either $|V|$ dense predicate evaluations or
one sparse predicate evaluation per candidate occurrence.  The abstract
balanced-reduction bound is $0$ for no messages and
$\lfloor\log_2 m\rfloor+1$ otherwise.

For comparison with an unfused pointwise schedule, let $M(Q)$ count map/zip
materialization passes and let $T_G(Q,C)$ count their full-stream temporary
value volume.  These are language-level quantities; they do not assert that a
particular compiler actually materializes or fuses a stream.

\begin{theorem}[Quantitative sharing-preserving normalization]
\label{thm:quantitative-normalization}
For every closed expression $Q$, graph $G$, and configuration $C$,
\[
\begin{aligned}
 |N(Q)| &= |Q|,&
 W_G(N(Q),C) &= W_G(Q,C),\\
 D_G(N(Q),C) &= D_G(Q,C),&
 L(N(Q)) &= L(Q),\\
 T_G(N(Q),C) &\leq T_G(Q,C),&
 M(N(Q)) &\leq M(Q).
\end{aligned}
\]
In particular, normalization has syntax-size factor exactly one and cannot
hide duplicated scalar work behind denotational equivalence.
\end{theorem}

\begin{proof}
Structural induction proves the four exact traversal equations
\path{Traversal.normalize_nodeCount}, \path{normalize_work},
\path{normalize_depth}, and \path{normalize_scalarTemporaries}; the map case
uses the additive definitions of \path{Term.let1}.  Constructor elimination
lifts them to \path{Program.normalize_nodeCount}, \path{normalize_work},
\path{normalize_depth}, and \path{normalize_scalarTemporaries}.  A normal form
has no mandatory intermediate pointwise edge column, yielding
\path{normalize_fullStreamTemporaryValues_le} and
\path{normalize_materializations_le}.
\end{proof}

Let $E$ encode a typed BSP program as a normal-form destination push and let
$D$ decode such a normal form as BSP.  These two auxiliary translations expose
the already verified nested-fold/flattened-fold connection.

\begin{theorem}[Normal-form syntactic equivalence]
\label{thm:syntax-equivalence}
For every typed BSP program $P$ and normal-form \TA{} program $Q$ over the
same signature,
\[
 D(E(P))=P,\qquad E(D(Q))=Q.
\]
Consequently $E$ is a bijection between BSP syntax and \TA{} normal forms.
\end{theorem}

\begin{proof}
Both equations follow by constructor elimination; no semantic equality is
used.  Lean checks them as \path{decode_encode} and \path{encode_decode} in
namespace \path{Typed.TraversalAlgebra}, with injectivity and surjectivity
recorded separately.
\end{proof}

For arbitrary compositional expressions define
\[
 \mathsf{toBSP}(Q)=D(N(Q)),\qquad
 \mathsf{ofBSP}(P)=R(E(P)).
\]
Because $R$ may introduce map nodes whose normalization adds explicit
bindings, the full-language round trips are intentionally semantic rather
than syntactic.  Lean records the BSP-side finite-run round trip as
\path{Program.toBSP_ofBSP_run_correct}; the normal-form translations $E$ and
$D$ themselves retain the exact syntactic bijection above.

\begin{theorem}[Monoidal Frontier BSP--\TA{} equivalence]
\label{thm:ta-bsp-equivalence}
For every fixed signature $\Sigma$, finite ordered graph $G$, initial barrier
configuration $C$, and $k\in\mathbb{N}$,
\[
 \mathsf{run}_{\mathrm{BSP}}(P,G,k,C)
 =\mathsf{run}_{\mathrm{TAExpr}}(\mathsf{ofBSP}(P),G,k,C)
\]
for every typed BSP program $P$, and
\[
 \mathsf{run}_{\mathrm{TAExpr}}(Q,G,k,C)
 =\mathsf{run}_{\mathrm{BSP}}(\mathsf{toBSP}(Q),G,k,C)
\]
for every closed compositional \TA{} expression $Q$.
Both equalities include the complete dense store and ordered frontier.
\end{theorem}

\begin{proof}
For normal forms, the \TA{} destination terminal folds the flattened sequence
$\trav_G(F)$.  Lemma~\ref{lem:nested-fold-flattening} identifies that fold
pointwise with the independently defined nested BSP fold; Lean checks the
central equation as \path{encode_inbox_correct} and the finite executions as
\path{encode_run_correct} and \path{decode_run_correct}.  Compose these results
with Theorem~\ref{thm:ta-normalization}.  Lean checks the two equations for the
full expression grammar as \path{Program.ofBSP_run_correct} and
\path{Program.toBSP_run_correct}.  Their existential forms,
\path{bsp_to_ta_complete} and \path{ta_to_bsp_complete}, state directly that
every program on either side has a semantics-equivalent program on the other.
\end{proof}

\begin{corollary}[Empty-frontier termination]
Under the conventional driver rule that execution halts when the frontier is
empty, a BSP program reaches a halting configuration if and only if its
compositional \TA{} embedding does; an arbitrary closed \TA{} expression
reaches one if and only if its normalized BSP translation does.
\end{corollary}

\begin{proof}
Theorem~\ref{thm:ta-bsp-equivalence} gives equality of the ordered frontier at
every step index, so existence of an index with the empty frontier is
preserved.  Lean checks both directions as \path{ofBSP_terminates_iff} and
\path{toBSP_terminates_iff}.
\end{proof}

\subsection{Sparse frontier order, multiplicity, and work}

Let $D_G^{\mathrm{nested}}(F)$ be the candidate stream produced directly by
the BSP source machine and let $D_G^{\mathrm{flat}}(F)$ project destinations
from $\trav_G(F)$.  Both are lists rather than sets: parallel edges, repeated
frontier entries, and repeated destinations remain observable.

\begin{theorem}[Sparse candidate preservation]
\label{thm:sparse-frontier}
For every graph and ordered frontier,
\[
 D_G^{\mathrm{flat}}(F)=D_G^{\mathrm{nested}}(F)
\]
as exact lists.  For every typed predicate $p$,
$\mathsf{sparsePreserve}(p)$ is a stable filter of this list.  It introduces no
occurrence; every accepted vertex retains exactly its candidate multiplicity,
and every rejected vertex has multiplicity zero.  If the candidate count is
$c$, its predicate work is exactly $cW(p)$; hence it is no greater than dense
selection's $|V|W(p)$ whenever $c\leq |V|$.
\end{theorem}

\begin{proof}
Induct on the frontier.  In the cons case, both constructions append the
destination projection of the active source row to their recursive tail; the
head projections are definitionally equal and the induction hypothesis closes
the tail.  Stable-filter order follows from the standard list-filter sublist
lemma.  For a fixed vertex, case analysis on the Boolean predicate gives the
original count in the accepted case and zero in the rejected case.  Expanding
the two work definitions gives $cW(p)$ and $|V|W(p)$, so monotonicity of
multiplication proves the conditional inequality.  Candidate duplicates can
make $c>|V|$, and the theorem deliberately does not erase that cost.  The Lean
declarations are mapped in Table~\ref{tab:correspondence}.
\end{proof}

Because both machines pass these equal candidate lists to the same frontier
policy, the step, finite-run, and equivalence theorems above cover both dense
and sparse policies without weakening ordered-frontier equality.

\subsection{Terminal observations}
\label{sec:terminal-observations}

Barrier configurations are not the right codomain for every terminal.  We
therefore use the disjoint observation type
\[
 \mathsf{Result}(A)=
   \mathsf{emitted}(A^*)
   +\mathsf{sourceReduced}(A^*)
   +\mathsf{destinationReduced}(V\to A).
\]
The distinct list constructors make the indexing contract explicit: emission
contains one value per traversed edge occurrence in traversal order, whereas
source reduction contains one value per frontier occurrence after folding its
adjacency row in order.  Destination reduction observes a dense store.

The TA observer evaluates an arbitrary compositional traversal expression and
invokes the corresponding public terminal.  Independently, the BSP observer
uses a typed edge term and nested frontier/adjacency recursion; it does not
call flattened traversal for its definition.

\begin{theorem}[Bidirectional terminal-observation equivalence]
\label{thm:terminal-observation}
Every TA emission, source-reduction, or destination-reduction observer has a
normalized BSP observer with exactly the same $\mathsf{Result}$.  Conversely,
every typed BSP observer has a reified TA observer with exactly the same
result.  The equalities preserve list order and all occurrence multiplicities.
Under the unit-cost model, TA-to-BSP normalization also preserves observation
work exactly for all three terminals.
\end{theorem}

\begin{proof}
Row induction proves ordered emission and source folds; frontier induction
proves their concatenation and every nested destination fold.  The generic
lemmas are \path{Observation.Proof.map_traverse}, \path{source_correct}, and
\path{destination_correct}.  Pointwise normalization and reification then
give \path{Terminal.toBSP_observe_correct} and
\path{BSP.Terminal.toTA_observe_correct}.  Exact work preservation is
\path{Observation.Cost.toBSP_work}.
\end{proof}

Thus semantic expressive completeness and the converse hold for the entire
stated closed expression grammar, not merely for a record chosen to be
isomorphic to BSP.  This result is relative to one arbitrary but fixed
signature $\Sigma$; it does not require translating $\Sigma$ to an execution
platform or proving an implementation.

\subsection{Orthogonal signature transport}

Signature transport is useful when changing scalar vocabularies, but is not a
premise of Theorem~\ref{thm:ta-bsp-equivalence}.

\begin{definition}[Signature lowering]
A lowering $L:\Sigma\Rightarrow\Sigma'$ maps every literal, primitive, and
monoid symbol to a symbol of the same object-language type in $\Sigma'$.  It
also carries the denotational equations
\[
 \denote{L(\ell)}_{\Sigma'}=
   \denote{\ell}_\Sigma,
 \qquad
 \denote{L(o)}_{\Sigma'}(x)=
   \denote{o}_\Sigma(x),
 \qquad
 \denote{L(M)}_{\Sigma'}=
   \denote{M}_\Sigma.
\]
The last equality is equality of the complete reduction operation, including
its identity and combine function.  Lowerings have an identity and compose by
composing symbol maps and their equations.
\end{definition}

Lowering extends structurally to terms, dense and sparse frontier policies,
and both source and target programs.  Variables, products, and let-binding
structure are unchanged.

\begin{lemma}[Term-lowering preservation]
\label{lem:term-lowering}
For every well-typed term $\Gamma\vdash_\Sigma t:\tau$, environment $\rho$,
and certified lowering $L:\Sigma\Rightarrow\Sigma'$,
\[
 \denote{\mathsf{lower}_L(t)}_{\Sigma'}\rho
 =\denote{t}_\Sigma\rho.
\]
\end{lemma}

\begin{proof}
Structural induction on $t$.  Variables are unchanged, literal and primitive
cases use the corresponding certificate equations, the product case uses two
induction hypotheses, and the let case preserves both the bound input and its
singleton body environment.  Lean checks this as
\path{Typed.Term.evaluate_lower}.  It also checks identity and composition for
terms, policies, and complete programs.
\end{proof}

\begin{theorem}[Compilation/lowering naturality]
\label{thm:lowering-naturality}
For every typed source program $P$ and certified lowering $L$,
\[
 \mathsf{compile}_{\mathrm{syn}}(\mathsf{lower}_L(P))
 =\mathsf{lower}_L(\mathsf{compile}_{\mathrm{syn}}(P)).
\]
Moreover, for every finite graph, initial configuration, and
$k\in\mathbb{N}$, the lowered source and the lowered compiled target produce
equal configurations after $k$ supersteps.
\end{theorem}

\begin{proof}
The syntactic square is fieldwise definitional equality and is checked by
\path{Typed.compile_lower}.  Lemma~\ref{lem:term-lowering} gives equal source
and target denotations before and after lowering; combining it with
Corollary~\ref{cor:typed-run} gives the finite-run square checked by
\path{Typed.lower_compile_run_correct}.
\end{proof}

No concrete execution platform occurs in these statements.  The non-identity
example in \path{TraversalAlgebra.TypedExample} merely witnesses transport
between two mathematical symbol datatypes.  It is not implementation
verification and is not needed by the BSP--\TA{} equivalence.

The proved boundary is now semantic and quantitative at the language level,
including all three terminal observations and both frontier policies.  The
remaining refinement programme is narrower:

\begin{enumerate}
  \item refine arbitrary parallel reduction-tree shapes to the
        permutation-invariant denotation already proved;
  \item replace unit scalar costs with signature- and backend-weighted costs,
        then refine transfer, allocation, residency, and CSR storage bounds;
  \item derive representative algorithms, including the Gunrock suite, as
        coverage evidence rather than as the completeness proof itself.
\end{enumerate}

The executable Massively/CubeCL realization and its property tests remain
artifact evidence for implementability, not an assumption of these formal
language results.  A Turing-machine simulation or a connection to fixed-point
logics could still be added as a secondary expressiveness result, but it would
say less about frontier-parallel structure than the correspondence proved
here.
